\chapter*{Abstract}
\makeatletter\@mkboth{Abstract}{Abstract}\makeatother
\addcontentsline{toc}{chapter}{Abstract}


In the last decade, machine learning has emerged as the driving force behind revolutionary advances in areas such as computer vision, natural language processing, and autonomous robotics. However, the success of the most powerful models, particularly in the supervised learning paradigm, critically depends on the availability of vast labeled datasets. The labeling process, often performed by human experts, constitutes a logistical and economic bottleneck, being in many cases a laborious, slow, and error-prone task. In this scenario, \acrfull{SSL} emerges as an essential discipline that seeks to maximize classifier performance using a small set of labeled data in combination with a large amount of unlabeled data.

This Bachelor's thesis addresses the problem of semi-supervised classification from an analytical perspective, analyzing how knowledge of the geometry of the input space allows refining decision boundaries. Unlike other more empirical approaches, this study focuses on the mathematical foundations that justify the use of unlabeled data. For these data to add value, certain regularity assumptions must be met. In this regard, the three fundamental assumptions of \acrshort{SSL} are formalized and discussed: the smoothness assumption, which postulates that nearby points in high-density regions should share the same label; the cluster or low-density assumption, which suggests that the decision boundary should pass through regions where data are scarce; and the manifold assumption, which assumes that high-dimensional data concentrate on a lower-dimensional structure. Understanding these assumptions is crucial, as if the data do not satisfy them, including unlabeled examples can introduce noise and degrade the performance of the original model.

The central contribution of this work is the proposal of a structured taxonomy to organize the vast literature of \acrshort{SSL}, inspired by recent surveys, especially \citet{van2020survey}, and adapted with motivated modifications to the aims of this study. Through a didactic yet rigorous approach, algorithms are classified according to their mathematical nature and how they incorporate unsupervised information. The work not only compiles scattered information but synthesizes it through a comparative analysis of the loss functions and regularization terms specific to each paradigm. The goal is to provide the reader with a conceptual map that allows discerning between inductive methods, designed to learn a generalizable decision rule for new data, and transductive methods, whose purpose is to infer labels exclusively for a previously defined set of unlabeled points.

The proposed taxonomy is developed in detail through four major blocks. First, wrapper methods, such as self-training and disagreement-based methods, are analyzed, which use supervised models iteratively to generate pseudo-labels. Second, unsupervised preprocessing is studied, where feature extraction and clustering techniques are employed to transform the representation space before applying the classifier. The third block, of particular mathematical interest, analyzes intrinsically semi-supervised methods. Here, generative models that use the EM algorithm to optimize mixtures of distributions and discriminative methods such as Semi-supervised Support Vector Machines (S3VM), which seek to maximize the margin in the presence of unlabeled data, are explored in depth. Finally, graph-based methods are explored, analyzing the construction of adjacency matrices and the use of harmonic functions and mincut for label propagation.

In conclusion, this work establishes a solid theoretical framework that clarifies the mathematical foundations of current \acrshort{SSL} techniques. The systematic review of algorithms and their properties highlights that the success of \acrshort{SSL} does not depend on the amount of data per se, but on the correct alignment between model assumptions and the underlying geometric structure of the data. This study lays the groundwork for future lines of research, such as the empirical evaluation of these taxonomies in large language models or advanced vision systems, ensuring a smooth transition between mathematical theory and practical implementation.